\input{../tex-inputs/latex-leseansicht-vorspann}

\section[Gustav Schwarzkopf an Arthur Schnitzler, 10. 8. 1911]{L04265 Gustav Schwarzkopf an Arthur Schnitzler, 10. 8. 1911}
\nopagebreak\mylabel{L04265v}
\rehead{ }\normalsize\beginnumbering\briefempfaengerindex{Schnitzler, Arthur@\textsc{Schnitzler, Arthur}!zzzSchwarzkopf, Gustav@\emph{von Gustav Schwarzkopf}!1911-08-101@{10. 8. 1911}|(be}
\toendnotes[C]{\smallbreak\pagebreak[2]}
\correspDesc{Versand  durch Gustav Schwarzkopf am 10. 8. 1911 in Wien
\newline{}Erhalt  durch Arthur Schnitzler im Zeitraum [11. 8. 1911 – 15. 8. 1911?] in Semmering}\toendnotes[C]{\smallbreak}
\Standort{CUL, Schnitzler, B 96a.}
\physDesc{Postkarte, 734 Zeichen
\newline{}Handschrift: schwarze Tinte, deutsche Kurrent
\newline{}Versand: Stempel: »\nobreak{}\oindex{I., Innere Stadt@\textbf{I., Innere Stadt}, \emph{Verwaltungsgebiet}|pwk}1/1 Wien 8, 10. VIII. 11, 8\nobreak{}«.  }\toendnotes[C]{\smallbreak}\pstart{}{\pb}I. Tiefer Graben 17.\oindex{Wien@\textbf{Wien}!I., Innere Stadt@\textbf{I., Innere Stadt}!Tiefer Graben 17@\textbf{Tiefer Graben 17}, \emph{Wohngebäude}|pw}\pend{}\pstart{}II. Th. 6.\pend{}{\bigskip}\pstart{}Herrn\pend{}\pstart{}\textsc{D\textsuperscript{r} Arthur
                  Schnitzler}\pend{}\pstart{}\textsc{Se{\geminationm}ering\oindex{Semmering@\textbf{Semmering}, \emph{Verwaltungsgebiet}|pw}}\pend{}\pstart{}\textsc{Südbahn-Hôtel\oindex{Südbahnhotel [Semmering]@\textbf{Südbahnhotel [Semmering]}, \emph{Hotel}|pw}}\pend{}{\bigskip}\vspace{1em}
\pstart
           \raggedleft{}{\pb}10. VIII. 11\pend
           \vspace{0.5em}
\pstart
           Lieber Arthur, ich war nach Mautern\oindex{Mautern in Steiermark@\textbf{Mautern in Steiermark}, \emph{Verwaltungsgebiet}|pw}
               noch 3 Tage in \textsc{Sekirn\oindex{Sekirn@\textbf{Sekirn}|pw}}, wo es beträchtlich warm war und bin{ }ſeit Ende Juli wieder hier.
               Vorigen So{\geminationn}tag Nachmittag habe ich bei Ihnen angefragt – ich dachte
               mir, vielleicht haben Sie Ihre Dispoſitionen geändert und{ }ſind zufällig doch noch in
                  Wien\oindex{Wien@\textbf{Wien}, \emph{Verwaltungsgebiet}|pw} – da erfuhr ich{ }ſchon von dem Stubenmädchen\pwindex{Loew, Anna *~11.\,4.\,1888 Ješín@\textsc{Loew, Anna} (*~11.\,4.\,1888 Ješín), \emph{Kinderbetreuerin, Dienstbotin}|pwuv}, daß Sie{ }ſich auf dem Se{\geminationm}ering\oindex{Semmering@\textbf{Semmering}, \emph{Verwaltungsgebiet}|pw} befinden. – Haben Sie vielleicht durch dritte Perſonen etwas
               von Richard\pwindex{Beer-Hofmann, Richard 11.\,7.\,1866 Wien – 26.\,9.\,1945 New York City@\textsc{Beer-Hofmann, Richard} (11.\,7.\,1866 Wien – 26.\,9.\,1945 New York City), \emph{Schriftsteller}|pw} gehört oder{ }ſollte er Ihnen gar – eine
               Anſichtskarte geſchickt haben? Ich weiß nichts von ihm. – Bitte, melden Sie Ihrem
               Herrn {\pb}Sohn\pwindex{Schnitzler, Heinrich 9.\,8.\,1902 Hinterbrühl – 12.\,7.\,1982 Wien@\textsc{Schnitzler, Heinrich} (9.\,8.\,1902 Hinterbrühl – 12.\,7.\,1982 Wien), \emph{Regisseur, Schauspieler}|pwv} meine verſpäteten Glückwünſche zu ſeinem neunten Geburtstage\pend
           
\pstart
           und grüßen Sie Ihre{\\[\baselineskip]}Frau\pwindex{Schnitzler, Olga 17.\,1.\,1882 Wien – 13.\,1.\,1970 Lugano@\textsc{Schnitzler, Olga} (17.\,1.\,1882 Wien – 13.\,1.\,1970 Lugano), \emph{Schauspielerin, Sängerin}|pwv} herzlichſt{\\[\baselineskip]}Ihr{\\[\baselineskip]}\spacefill\mbox{Gustav.}\pend
           \leftskip=0em{}\selectlanguage{ngerman}\endnumbering\briefempfaengerindex{Schnitzler, Arthur@\textsc{Schnitzler, Arthur}!zzzSchwarzkopf, Gustav@\emph{von Gustav Schwarzkopf}!1911-08-101@{10. 8. 1911}|)be}\mylabel{L04265h}
\begin{anhang}
\end{anhang}\newcommand{\dateiname}{L04265}\newcommand{\titel}{Gustav Schwarzkopf an Arthur Schnitzler, 10. 8. 1911}\newcommand{\editorInnen}{Herausgegeben von Jahnke, SelmaMüller, Martin Anton}\input{../tex-inputs/latex-leseansicht-abspann}
